% User manual of the RAG on Medical Documents app, for a reader who is not a developer.
% Build from this folder: pdflatex main.tex
\documentclass[11pt,a4paper]{article}

\usepackage[T1]{fontenc}
\usepackage[utf8]{inputenc}
\usepackage{lmodern}
\usepackage{microtype}
\usepackage[a4paper,margin=2.3cm,headheight=14pt]{geometry}
\usepackage{booktabs}
\usepackage{tabularx}
\usepackage{array}
\usepackage{xcolor}
\usepackage{enumitem}
\usepackage{fancyhdr}
\usepackage{caption}
\usepackage{float}
\usepackage{listings}
\usepackage{xurl}
\usepackage[hidelinks]{hyperref}

\newcommand{\TODO}[1]{\textcolor{red}{\textbf{[TODO: #1]}}}
\newcommand{\code}[1]{\texttt{\detokenize{#1}}}
\newcommand{\ui}[1]{\textsf{\textbf{#1}}}  % a label shown in the app
\newcommand{\shot}[2]{%
  \begin{figure}[H]\centering
  \fbox{\parbox[c][#1][c]{0.85\textwidth}{\centering\TODO{screenshot: #2}}}
  \end{figure}}

\lstset{basicstyle=\ttfamily\small, breaklines=true, columns=fullflexible,
        frame=single, rulecolor=\color{black!25}, backgroundcolor=\color{black!3},
        xleftmargin=0.5em, xrightmargin=0.5em, aboveskip=0.5em, belowskip=0.5em}
\setlist{itemsep=0.2em, topsep=0.3em}
\setlength{\parskip}{0.4em}
\setlength{\parindent}{0pt}
\renewcommand{\arraystretch}{1.2}

\pagestyle{fancy}
\fancyhf{}
\fancyhead[L]{\small \TODO{school name}}
\fancyhead[R]{\small Group \TODO{NN} -- RAG on Medical Documents}
\fancyfoot[C]{\small \thepage}

\begin{document}

\begin{center}
  {\LARGE\bfseries RAG on Medical Documents}\\[0.3em]
  {\Large User manual}\\[0.6em]
  Boussoualim Yassine \quad Benahcene Mohamed Wassim \quad Aissaoui Anes Safouane \quad Djerdali Mohamed Achraf\\[0.2em]
  Group \TODO{group number} \quad -- \quad \TODO{school header}\\
  LLM Final Projects 2026-2027, Prof.\ A.\ Moussaoui, Project 2
\end{center}

\section*{What this tool does}

You type a medical question. The tool looks for the passages that answer it in 50 documents: 20
guidelines of the World Health Organization (WHO) and 30 open-access articles from PubMed Central.
A free language model then writes an answer from those passages only, and names its source after
each claim. If the documents do not contain the answer, the tool says so instead of guessing:

\begin{center}
\emph{I cannot answer this question from the provided documents.}
\end{center}

\noindent\fcolorbox{red!60!black}{red!4}{\parbox{0.96\textwidth}{\textbf{Important.}
This tool is a student project. It is not medical advice. Check every answer with a clinician.
Questions are sent to OpenRouter, an online service: do not type anything that identifies a
patient (name, date of birth, file number, address).}}

The app shows this warning at the top of the page and under every answer.

\section*{1. What you need}

\begin{itemize}[leftmargin=1.5em]
  \item A Linux computer with Python 3.12 and an internet connection.
  \item Several gigabytes of free disk space: the tool downloads its own language models (the largest is about 2\,GB) and the documents.
  \item A free key from OpenRouter (\url{https://openrouter.ai}). The key is a long text that lets
        the tool use the free models of your account. The tool only uses free models.
  \item Time: building the search index takes about 3 hours on a laptop, but only once. After that,
        the app starts much faster.
\end{itemize}

All the commands below are typed in a terminal, in the project folder (the folder that contains
\code{README.md}). Type each line, then press Enter.

\section*{2. Install}

\begin{lstlisting}
python3 -m venv .venv
.venv/bin/pip install -r requirements.txt
\end{lstlisting}

The first line creates a private Python environment in the folder \code{.venv}. The second
installs the libraries the tool needs, including PyTorch; the download can take several
minutes. On Debian or Ubuntu, if the first line fails, install the missing package first with
\code{sudo apt install python3-venv}, then try again.

\section*{3. Put your key in \texttt{.env}}

\begin{lstlisting}
cp .env.example .env
\end{lstlisting}

Open the new file \code{.env} with a text editor. It has two lines. Paste your key right after the
equals sign of the first line, with no space and no quotes:

\begin{lstlisting}
OPENROUTER_API_KEY=your-key-here
OPENROUTER_MODEL=nvidia/nemotron-3-ultra-550b-a55b:free
\end{lstlisting}

Save the file. Leave the second line as it is. Keep this file private: never send it, never put it
online, never add it to git.

\section*{4. Get the documents and build the index}

\begin{lstlisting}
.venv/bin/python scripts/fetch_corpus.py --from-manifest
.venv/bin/python src/index.py
\end{lstlisting}

The first command downloads the 50 documents. It prints one line per file and ends with a count
such as \code{done: 50/50 files written}. If a few files failed, run it again.

The second command builds the search index: it reads every document, cuts it into short passages
and computes a numeric fingerprint of each one. The first time, it downloads a model of about 2\,GB.
Then it prints its progress regularly, with an estimate of the minutes left. It takes about
3 hours on a laptop. You can stop it (Ctrl+C) or lose power: run the same command again and it
continues where it stopped. It ends with a line that starts with \code{indexed}. You only do this
once.

\section*{5. Start the app}

\begin{lstlisting}
.venv/bin/python src/app.py
\end{lstlisting}

The first start downloads one more model (about 90\,MB). When the terminal shows the address, open
\url{http://127.0.0.1:7860} in your web browser. The app runs on your own computer; leave the
terminal open while you use it, and press Ctrl+C in the terminal to stop it.

\shot{4.2cm}{the app's page when it opens, empty (\code{07_EXTRA/screenshots/})}

\section*{6. Ask a question}

\begin{enumerate}[leftmargin=1.5em]
  \item Type your question in the \ui{Question} box, in English, as one full sentence. Or click one
        of the three examples under the page; they fill the box for you.
  \item Click \ui{Ask}, or press Enter in the box.
  \item Wait. An answer usually takes a few seconds; in our tests between about 2 and 18 seconds.
        A refusal often comes back in under a second, because the model is not asked at all.
\end{enumerate}

The other controls can stay as they are:

\begin{itemize}[leftmargin=1.5em]
  \item \ui{Model}: which free language model writes the answer. The default is
        \nolinkurl{nvidia/nemotron-3-ultra-550b-a55b:free}. If one model keeps failing, try
        another. (The app always starts with this default, whatever \code{.env} says.)
  \item \ui{Retrieval and refusal gate}: how passages are found and how the tool decides to refuse.
        Keep the default, \ui{Hybrid + cross-encoder gate}. \ui{Dense} and
        \ui{Hybrid (dense + BM25)} are simpler methods, kept for comparison.
  \item \ui{Refusal threshold}: how sure the tool must be before it answers. Higher means more
        refusals. Changing the mode puts the slider back to that mode's default (0.5 for the
        cross-encoder gate, 0.6 for the two others).
\end{itemize}

\section*{7. Read the answer and its citations}

Each claim of an answer is followed by its source in square brackets:

\begin{itemize}[leftmargin=1.5em]
  \item A WHO guideline, with the file name and the page number, for example\\
        \code{[WHO_09_Guidelines_for_the_treatment_of_malaria.pdf, p. 14]}.\\
        Open that PDF at that page to check the claim.
  \item \code{[article title]}: a PubMed Central article, named by its title. Articles have no page.
\end{itemize}

Under the answer, a table lists the five passages the tool found, best first:

\begin{center}\small
\begin{tabularx}{0.95\textwidth}{@{}l X@{}}
  \toprule
  Column & Meaning \\
  \midrule
  \ui{\#} & Rank of the passage, 1 to 5. \\
  \ui{Source} & File name of the WHO guideline, or title of the article. \\
  \ui{Page} & Page in the PDF; \code{-} for an article. \\
  \ui{Similarity} & How close the passage is to the question, from 0 to 1; \code{-} for a passage found only by keyword search. \\
  \ui{Cited} & \code{yes} if the answer cites this passage. \\
  \bottomrule
\end{tabularx}
\end{center}

The tool checks that every citation points at one of these five passages. It does \emph{not}
check that the passage really says what the sentence claims. Always read the cited page before you
rely on an answer.

\shot{4.6cm}{an answer to ``How long should parenteral artesunate be given to a patient with severe
malaria?'' with its sources table (\code{07_EXTRA/screenshots/})}

\section*{8. What a refusal means}

A refusal is the sentence \emph{I cannot answer this question from the provided documents.},
followed by a short reason. It is not an error. It means the tool did not find support for an
answer in its 50 documents, so it did not guess. The reason tells you why:

\begin{center}\small
\begin{tabularx}{0.95\textwidth}{@{}>{\raggedright\arraybackslash}p{6.2cm} X@{}}
  \toprule
  Reason shown & What it means \\
  \midrule
  The cross-encoder rates the best passage 0.16, below the threshold. & None of the five passages answers the question (default mode). \\
  Best similarity 0.41 is below the threshold. & No passage is close enough to the question (other modes). \\
  The retrieved passages do not contain the answer. & The model read the passages and found no answer in them. \\
  The model's answer did not cite the retrieved passages. & The model wrote an answer without valid citations, so it was not shown. \\
  \bottomrule
\end{tabularx}
\end{center}

The numbers change with each question. The documents cover hypertension, diabetes, malaria,
tuberculosis, HIV and antenatal care, among others; many questions outside them are refused. You
can rephrase the question. Do not lower the threshold just to get an answer: a lower threshold
lets in passages that do not answer the question.

\shot{3.8cm}{a refusal, for example ``What is the first-line topical treatment for scabies?''
(\code{07_EXTRA/screenshots/})}

\section*{9. Troubleshooting}

\begin{center}\small
\begin{tabularx}{\textwidth}{@{}>{\raggedright\arraybackslash}p{6.6cm} X@{}}
  \toprule
  What you see & What to do \\
  \midrule
  In the terminal, the app does not start:\newline
  \code{OPENROUTER_API_KEY is missing: copy .env.example to .env and add your key.}
  & Do step 3: create \code{.env} and paste your key after the equals sign. \\
  In the terminal, the app does not start:\newline
  \code{The index is empty: run .venv/bin/python src/index.py first.}
  & Do step 4. If the build was interrupted, run the same command again to finish it. \\
  In the app: \code{Type a question first.} & The question box was empty. \\
  In the app: \code{The model call failed: ...} followed by \code{HTTP Error 401} & The key in \code{.env} is wrong or was revoked. Paste a valid key, then restart the app. \\
  In the app: \code{The model call failed: the daily free-model limit of the OpenRouter account is used up; it resets at 00:00 UTC} & The account's free requests for today are used up (50 a day). Try again after 00:00 UTC (01:00 in Paris in winter). \\
  In the app: \code{The model call failed: ...} with another message, after a long wait & The free model is busy. The tool already retried several times. Try again later or pick another \ui{Model}. \\
  The browser cannot open \url{http://127.0.0.1:7860} & Check that the terminal running the app is still open and shows no error. \\
  \code{python3 -m venv .venv} fails & On Debian or Ubuntu: \code{sudo apt install python3-venv}, then repeat step 2. \\
  \bottomrule
\end{tabularx}
\end{center}

\vfill
{\small\textit{Claude (Anthropic) helped draft this manual.}}

\end{document}
